\documentclass[11pt]{article}
\usepackage{amssymb}
\usepackage{amsthm}
\usepackage{enumitem}
\usepackage{physics,amsmath}
\usepackage{bm}
\usepackage{adjustbox}
\usepackage{mathrsfs}
\usepackage{graphicx}
\usepackage{siunitx}
\usepackage[mathscr]{euscript}


\title{\textbf{Solved selected problems of General Relativity - Thomas A. Moore}}
\author{Franco Zacco}
\date{}

\addtolength{\topmargin}{-3cm}
\addtolength{\textheight}{3cm}

\newcommand{\hatr}{\bm{\hat{r}}}
\newcommand{\hatn}{\bm{\hat{n}}}
\newcommand{\hatx}{\bm{\hat{x}}}
\newcommand{\haty}{\bm{\hat{y}}}
\newcommand{\hatz}{\bm{\hat{z}}}
\newcommand{\hatth}{\bm{\hat{\theta}}}
\newcommand{\hatphi}{\bm{\hat{\phi}}}
\newcommand{\hatrho}{\bm{\hat{\rho}}}
\newcommand{\esi}[1]{\bm{e}_{#1}}
\newcommand{\etht}{\bm{e}_\theta}
\newcommand{\sci}[1]{\times 10^{#1}}

\theoremstyle{definition}
\newtheorem*{solution*}{Solution}
\renewcommand*{\proofname}{Solution}

\begin{document}
\maketitle
\thispagestyle{empty}

\section*{Chapter 19 - The Riemann Tensor}

\begin{proof}{\textbf{BOX 19.1} - Exercise 19.1.1.}\\
We know that in a LIF
\begin{align*}
    R_{\alpha\beta\mu\nu} = g_{\alpha\sigma}R^{\sigma}_{\beta\mu\nu}
    = g_{\alpha\sigma}[\partial_{\mu}\Gamma_{\beta\nu}^\sigma
    - \partial_{\nu}\Gamma_{\beta\mu}^\sigma]
\end{align*}
Replacing the value of $\Gamma_{\beta\nu}^\sigma$ and $\Gamma_{\beta\mu}^\sigma$
we get that
\begin{align*}
    R_{\alpha\beta\mu\nu}
    &= g_{\alpha\sigma}[\partial_{\mu}\Gamma_{\beta\nu}^\sigma
    - \partial_{\nu}\Gamma_{\beta\mu}^\sigma]\\
    %
    &= \frac{1}{2}g_{\alpha\sigma}[
    \partial_{\mu}(g^{\sigma\gamma}[\partial_{\beta}g_{\nu\gamma}
    + \partial_{\nu}g_{\gamma\beta} - \partial_{\gamma}g_{\beta\nu}])\\
    &\quad
    - \partial_{\nu}(g^{\sigma\gamma}[\partial_{\beta}g_{\mu\gamma}
    + \partial_{\mu}g_{\gamma\beta} - \partial_{\gamma}g_{\beta\mu}])]\\
    %
    &= \frac{1}{2}g_{\alpha\sigma}[
    \partial_{\mu}g^{\sigma\gamma}[\partial_{\beta}g_{\nu\gamma}
    + \partial_{\nu}g_{\gamma\beta} - \partial_{\gamma}g_{\beta\nu}]
    + g^{\sigma\gamma}[\partial_{\mu}\partial_{\beta}g_{\nu\gamma}
    + \partial_{\mu}\partial_{\nu}g_{\gamma\beta}
    - \partial_{\mu}\partial_{\gamma}g_{\beta\nu}]\\
    &\quad
    - \partial_{\nu}g^{\sigma\gamma}[\partial_{\beta}g_{\mu\gamma}
    + \partial_{\mu}g_{\gamma\beta} - \partial_{\gamma}g_{\beta\mu}]
    - g^{\sigma\gamma}[\partial_{\nu}\partial_{\beta}g_{\mu\gamma}
    + \partial_{\nu}\partial_{\mu}g_{\gamma\beta}
    - \partial_{\nu}\partial_{\gamma}g_{\beta\mu}]]\\
    %
    &= \frac{1}{2}(g_{\alpha\sigma}
    g^{\sigma\gamma}
    [\partial_{\mu}\partial_{\beta}g_{\nu\gamma}
    + \partial_{\mu}\partial_{\nu}g_{\gamma\beta}
    - \partial_{\mu}\partial_{\gamma}g_{\beta\nu}]\\
    &\quad
    - g_{\alpha\sigma}g^{\sigma\gamma}
    [\partial_{\nu}\partial_{\beta}g_{\mu\gamma}
    + \partial_{\nu}\partial_{\mu}g_{\gamma\beta}
    - \partial_{\nu}\partial_{\gamma}g_{\beta\mu}])\\
    %
    &= \frac{1}{2}(
    \delta^\gamma_{\alpha}[\partial_{\mu}\partial_{\beta}g_{\nu\gamma}
    + \partial_{\mu}\partial_{\nu}g_{\gamma\beta}
    - \partial_{\mu}\partial_{\gamma}g_{\beta\nu}]\\
    &\quad
    - \delta^\gamma_{\alpha}[\partial_{\nu}\partial_{\beta}g_{\mu\gamma}
    + \partial_{\nu}\partial_{\mu}g_{\gamma\beta}
    - \partial_{\nu}\partial_{\gamma}g_{\beta\mu}])\\
    %
    &= \frac{1}{2}(
    \partial_{\mu}\partial_{\beta}g_{\nu\alpha}
    + \partial_{\mu}\partial_{\nu}g_{\alpha\beta}
    - \partial_{\mu}\partial_{\alpha}g_{\beta\nu}
    - \partial_{\nu}\partial_{\beta}g_{\mu\alpha}
    - \partial_{\nu}\partial_{\mu}g_{\alpha\beta}
    + \partial_{\nu}\partial_{\alpha}g_{\beta\mu})\\
    %
    &= \frac{1}{2}(
    \partial_{\beta}\partial_{\mu}g_{\alpha\nu}
    + \partial_{\alpha}\partial_{\nu}g_{\beta\mu}
    - \partial_{\beta}\partial_{\nu}g_{\alpha\mu}
    - \partial_{\alpha}\partial_{\mu}g_{\beta\nu})
\end{align*}
In the last step, we used the symmetry of the metric and that the order of
the partial derivatives doesn't matter.
\end{proof}

\cleardoublepage
\begin{proof}{\textbf{BOX 19.2} - Exercise 19.2.1.}\\
We know that
\begin{align*}
    R_{\alpha\beta\mu\nu} &= \frac{1}{2}(
    \partial_{\beta}\partial_{\mu}g_{\alpha\nu}
    + \partial_{\alpha}\partial_{\nu}g_{\beta\mu}
    - \partial_{\beta}\partial_{\nu}g_{\alpha\mu}
    - \partial_{\alpha}\partial_{\mu}g_{\beta\nu})
\end{align*}
So 
\begin{align*}
    -R_{\alpha\beta\mu\nu} &= \frac{1}{2}(
    \partial_{\alpha}\partial_{\mu}g_{\beta\nu}
    + \partial_{\beta}\partial_{\nu}g_{\alpha\mu}
    - \partial_{\alpha}\partial_{\nu}g_{\beta\mu}
    - \partial_{\beta}\partial_{\mu}g_{\alpha\nu}
    )
\end{align*}
But what we have on the LHS is just $R_{\alpha\beta\mu\nu}$ but with $\alpha$
and $\beta$ swapped, hence
\begin{align*}
    -R_{\alpha\beta\mu\nu} &= \frac{1}{2}(
    \partial_{\alpha}\partial_{\mu}g_{\beta\nu}
    + \partial_{\beta}\partial_{\nu}g_{\alpha\mu}
    - \partial_{\alpha}\partial_{\nu}g_{\beta\mu}
    - \partial_{\beta}\partial_{\mu}g_{\alpha\nu}
    ) = R_{\beta\alpha\mu\nu}
\end{align*}
\end{proof}
\begin{proof}{\textbf{BOX 19.2} - Exercise 19.2.2.}\\
Taking
\begin{align*}
    R_{\alpha\beta\mu\nu} &= \frac{1}{2}(
    \partial_{\beta}\partial_{\mu}g_{\alpha\nu}
    + \partial_{\alpha}\partial_{\nu}g_{\beta\mu}
    - \partial_{\beta}\partial_{\nu}g_{\alpha\mu}
    - \partial_{\alpha}\partial_{\mu}g_{\beta\nu})
\end{align*}
But renaming the indices as $\alpha\to \mu$, $\beta\to \nu$, $\mu\to \alpha$
and $\nu\to \beta$ we get that
\begin{align*}
    R_{\mu\nu\alpha\beta} 
    &= \frac{1}{2}(
    \partial_{\nu}\partial_{\alpha}g_{\mu\beta}
    + \partial_{\mu}\partial_{\beta}g_{\nu\alpha}
    - \partial_{\nu}\partial_{\beta}g_{\mu\alpha}
    - \partial_{\mu}\partial_{\alpha}g_{\nu\beta})\\
    &= \frac{1}{2}(
    \partial_{\beta}\partial_{\mu}g_{\alpha\nu}
    + \partial_{\alpha}\partial_{\nu}g_{\beta\mu}
    - \partial_{\beta}\partial_{\nu}g_{\alpha\mu}
    - \partial_{\alpha}\partial_{\mu}g_{\beta\nu})
\end{align*}
Where In the last step, we used the symmetry of the metric and that the order
of the partial derivatives doesn't matter.
\\
Therefore
\begin{align*}
    R_{\mu\nu\alpha\beta} = R_{\alpha\beta\mu\nu}
\end{align*}
\end{proof}

\cleardoublepage
\begin{proof}{\textbf{BOX 19.2} - Exercise 19.2.3.}\\
We note that
\begin{align*}
    &R_{\alpha\beta\mu\nu} + R_{\alpha\nu\beta\mu} + R_{\alpha\mu\nu\beta} =\\
    &=\frac{1}{2}(
    \partial_{\beta}\partial_{\mu}g_{\alpha\nu}
    + \partial_{\alpha}\partial_{\nu}g_{\beta\mu}
    - \partial_{\beta}\partial_{\nu}g_{\alpha\mu}
    - \partial_{\alpha}\partial_{\mu}g_{\beta\nu})\\
    &\quad + \frac{1}{2}(
    \partial_{\nu}\partial_{\beta}g_{\alpha\mu}
    + \partial_{\alpha}\partial_{\mu}g_{\nu\beta}
    - \partial_{\nu}\partial_{\mu}g_{\alpha\beta}
    - \partial_{\alpha}\partial_{\beta}g_{\nu\mu})\\
    &\quad + \frac{1}{2}(
    \partial_{\mu}\partial_{\nu}g_{\alpha\beta}
    + \partial_{\alpha}\partial_{\beta}g_{\mu\nu}
    - \partial_{\mu}\partial_{\beta}g_{\alpha\nu}
    - \partial_{\alpha}\partial_{\nu}g_{\mu\beta})\\
    %
    &=\frac{1}{2}(
    (\partial_{\beta}\partial_{\mu}g_{\alpha\nu}
    - \partial_{\mu}\partial_{\beta}g_{\alpha\nu})
    + (\partial_{\alpha}\partial_{\nu}g_{\beta\mu}
    - \partial_{\alpha}\partial_{\nu}g_{\mu\beta})\\
    &\quad
    + (\partial_{\nu}\partial_{\beta}g_{\alpha\mu}
    - \partial_{\beta}\partial_{\nu}g_{\alpha\mu})
    + (\partial_{\alpha}\partial_{\mu}g_{\nu\beta}
    - \partial_{\alpha}\partial_{\mu}g_{\beta\nu})\\
    &\quad + 
    (\partial_{\mu}\partial_{\nu}g_{\alpha\beta}
    - \partial_{\nu}\partial_{\mu}g_{\alpha\beta})
    + (\partial_{\alpha}\partial_{\beta}g_{\mu\nu}
    - \partial_{\alpha}\partial_{\beta}g_{\nu\mu}))\\
    %
    &= 0
\end{align*}

\end{proof}

\cleardoublepage
\begin{proof}{\textbf{BOX 19.3} - Exercise 19.3.1.}\\
Given that $R_{\mu\nu\alpha\beta} = R_{\alpha\beta\mu\nu}$ then for example,
the component $R_{0201}$ is equal to $R_{0102}$, $R_{0301} = R_{0103}$ and
so on, so from the matrix we have in (19.16) only the upper (or lower) 
elements plus the diagonal elements are independent, then we have
$6 + 5+ 4+ 3+ 2+ 1 = 21$ independent components.
\end{proof}
\begin{proof}{\textbf{BOX 19.3} - Exercise 19.3.2.}\\
Let $\alpha = \beta = x$ where $x \in \{0,1,2,3\}$ then the equation (19.4d)
becomes
\begin{align*}
    R_{xx\mu\nu} + R_{x\nu x\mu} + R_{x\mu\nu x}
    &= 0 + R_{x\nu x\mu} - R_{x\mu x\nu}
    = R_{x\nu x\mu} - R_{x\nu x\mu}
    = 0
\end{align*}
\end{proof}
\begin{proof}{\textbf{BOX 19.3} - Exercise 19.3.3.}\\
In a two-dimensional space we have that $\alpha, \beta, \mu, \nu \in \{0,1\}$
and we know that the Riemann tensor is nonzero if $\alpha \neq\beta$ and
$\mu\neq\nu$, so the only possible components are
\begin{align*}
    R_{0101} \quad R_{0110} \quad R_{1001} \quad R_{1010}
\end{align*}
But also we know that $R_{\alpha\beta\mu\nu} = -R_{\alpha\beta\nu\mu}$
and $R_{\alpha\beta\mu\nu} = -R_{\beta\alpha\mu\nu}$ then
\begin{align*}
    R_{0101} = -R_{0110} = R_{1010} = -R_{1001}
\end{align*}
Therefore there is only one independent component of the Riemann tensor.
\end{proof}

\cleardoublepage
\begin{proof}{\textbf{BOX 19.4} - Exercise 19.4.1.}\\
From BOX 19.1 we know that 
\begin{align*}
    R_{\alpha\beta\mu\nu} &= \frac{1}{2}(
    \partial_{\beta}\partial_{\mu}g_{\alpha\nu}
    + \partial_{\alpha}\partial_{\nu}g_{\beta\mu}
    - \partial_{\beta}\partial_{\nu}g_{\alpha\mu}
    - \partial_{\alpha}\partial_{\mu}g_{\beta\nu})
\end{align*}
Changing $\mu\to\sigma$ and $\nu\to\mu$ we get that
\begin{align*}
    R_{\alpha\beta\sigma\mu} &= \frac{1}{2}(
        \partial_{\beta}\partial_{\sigma}g_{\alpha\mu}
        + \partial_{\alpha}\partial_{\mu}g_{\beta\sigma}
        - \partial_{\beta}\partial_{\mu}g_{\alpha\sigma}
        - \partial_{\alpha}\partial_{\sigma}g_{\beta\mu})
\end{align*}
Or changing $\mu\to\nu$ and $\nu\to\sigma$ we get that
\begin{align*}
    R_{\alpha\beta\nu\sigma} &= \frac{1}{2}(
    \partial_{\beta}\partial_{\nu}g_{\alpha\sigma}
    + \partial_{\alpha}\partial_{\sigma}g_{\beta\nu}
    - \partial_{\beta}\partial_{\sigma}g_{\alpha\nu}
    - \partial_{\alpha}\partial_{\nu}g_{\beta\sigma})
\end{align*}
Then
\begin{align*}
    &\partial_{\sigma} R_{\alpha\beta\mu\nu}
    + \partial_{\nu} R_{\alpha\beta\sigma\mu} 
    + \partial_{\mu} R_{\alpha\beta\nu\sigma} =\\
    &=\frac{1}{2}(
    \partial_{\sigma}\partial_{\beta}\partial_{\mu}g_{\alpha\nu}
    + \partial_{\sigma}\partial_{\alpha}\partial_{\nu}g_{\beta\mu}
    - \partial_{\sigma}\partial_{\beta}\partial_{\nu}g_{\alpha\mu}
    - \partial_{\sigma}\partial_{\alpha}\partial_{\mu}g_{\beta\nu})\\
    &\quad + \frac{1}{2}(
    \partial_{\nu}\partial_{\beta}\partial_{\sigma}g_{\alpha\mu}
    + \partial_{\nu}\partial_{\alpha}\partial_{\mu}g_{\beta\sigma}
    - \partial_{\nu}\partial_{\beta}\partial_{\mu}g_{\alpha\sigma}
    - \partial_{\nu}\partial_{\alpha}\partial_{\sigma}g_{\beta\mu})\\
    &\quad + \frac{1}{2}(
    \partial_{\mu}\partial_{\beta}\partial_{\nu}g_{\alpha\sigma}
    + \partial_{\mu}\partial_{\alpha}\partial_{\sigma}g_{\beta\nu}
    - \partial_{\mu}\partial_{\beta}\partial_{\sigma}g_{\alpha\nu}
    - \partial_{\mu}\partial_{\alpha}\partial_{\nu}g_{\beta\sigma})\\
    %
    &=\frac{1}{2}[
    (\partial_{\sigma}\partial_{\beta}\partial_{\mu}g_{\alpha\nu}
    - \partial_{\mu}\partial_{\beta}\partial_{\sigma}g_{\alpha\nu})
    + (\partial_{\sigma}\partial_{\alpha}\partial_{\nu}g_{\beta\mu}
    - \partial_{\nu}\partial_{\alpha}\partial_{\sigma}g_{\beta\mu})\\
    &\quad
    + (\partial_{\nu}\partial_{\beta}\partial_{\sigma}g_{\alpha\mu}
    - \partial_{\sigma}\partial_{\beta}\partial_{\nu}g_{\alpha\mu})
    + (\partial_{\nu}\partial_{\alpha}\partial_{\mu}g_{\beta\sigma}
    - \partial_{\mu}\partial_{\alpha}\partial_{\nu}g_{\beta\sigma})\\
    &\quad
    +(\partial_{\mu}\partial_{\beta}\partial_{\nu}g_{\alpha\sigma}
    - \partial_{\nu}\partial_{\beta}\partial_{\mu}g_{\alpha\sigma})
    + (\partial_{\mu}\partial_{\alpha}\partial_{\sigma}g_{\beta\nu}
    - \partial_{\sigma}\partial_{\alpha}\partial_{\mu}g_{\beta\nu}) ]\\
    %
    &= 0
\end{align*}
Since this should be valid for all frames then must be that
\begin{align*}
    \nabla_{\sigma} R_{\alpha\beta\mu\nu}
    + \nabla_{\nu} R_{\alpha\beta\sigma\mu} 
    + \nabla_{\mu} R_{\alpha\beta\nu\sigma} = 0
\end{align*}
\end{proof}

\cleardoublepage
\begin{proof}{\textbf{BOX 19.5} - Exercise 19.5.1.}\\
The Ricci tensor is defined as 
\begin{align*}
    R_{\beta\nu} = R^{\alpha}_{\beta\alpha\nu} = g^{\alpha\mu} R_{\alpha\beta\mu\nu}
\end{align*}
Let us consider $R_{\nu\beta}$, we see that
\begin{align*}
    R_{\nu\beta} &= g^{\alpha\mu} R_{\alpha\nu\mu\beta}
    = g^{\alpha\mu} R_{\mu\beta\alpha\nu}
    = g^{\mu\alpha} R_{\mu\beta\alpha\nu}
\end{align*}
Where we used that $R_{\alpha\nu\mu\beta} = R_{\mu\beta\alpha\nu}$ and that
$g^{\alpha\mu} = g^{\mu\alpha}$.
\\
So if we rename indexes $\mu \to \alpha$ and $\alpha \to \mu$ we get that
\begin{align*}
    R_{\nu\beta} = g^{\alpha\mu} R_{\alpha\beta\mu\nu} = R_{\beta\mu}
\end{align*}
Therefore, the Ricci tensor is symmetric.
\end{proof}

\cleardoublepage
\begin{proof}{\textbf{BOX 19.6} - Exercise 19.6.1.}\\
Let $\mu = \theta$ then the geodesic equation is
\begin{align*}
    \dv{\tau}(g_{\theta \nu}\dv{x^\nu}{\tau})
    - \frac{1}{2}\partial_\theta g_{\alpha\beta}\dv{x^\alpha}{\tau}\dv{x^\beta}{\tau}
    &= 0\\
    \dv{\tau}(g_{\theta \theta}\dv{\theta}{\tau})
    - \frac{1}{2}\partial_\theta g_{\phi\phi}\dv{\phi}{\tau}\dv{\phi}{\tau}
    &= 0\\
    \dv{\tau}(r^2\dv{\theta}{\tau})
    - r^2\sin\theta\cos\theta\dv{\phi}{\tau}\dv{\phi}{\tau}
    &= 0\\
    r^2\dv[2]{\theta}{\tau}
    - r^2\sin\theta\cos\theta\dv{\phi}{\tau}\dv{\phi}{\tau}
    &= 0\\
    \dv[2]{\theta}{\tau}
    - \sin\theta\cos\theta\dv{\phi}{\tau}\dv{\phi}{\tau}
    &= 0
\end{align*}
Hence the only nonzero Christoffel symbol with $\theta$ as a superscript is
\begin{align*}
    \Gamma^\theta_{\phi\phi} = -\sin\theta\cos\theta
\end{align*}
\end{proof}
\begin{proof}{\textbf{BOX 19.6} - Exercise 19.6.2.}\\
Let us compute $R_{\theta\phi\theta\phi}$ as follows
\begin{align*}
    R_{\theta\phi\theta\phi}
    &= r^2(\partial_{\theta}\Gamma^\theta_{\phi\phi} - \Gamma^\theta_{\phi\phi}
    \Gamma^\phi_{\phi\theta})\\
    &= r^2((\sin^2\theta - \cos^2\theta) + \sin\theta\cos\theta\cdot \cot\theta)\\
    &= r^2(\sin^2\theta - \cos^2\theta
    + \sin\theta\cos\theta\cdot \frac{\cos\theta}{\sin\theta})\\
    &= r^2(\sin^2\theta - \cos^2\theta + \cos^2\theta)\\
    &= r^2\sin^2\theta
\end{align*}
\end{proof}
\begin{proof}{\textbf{BOX 19.6} - Exercise 19.6.3.}\\
The only nonzero components of the Ricci tensor are $R_{\theta\theta}$ and
$R_{\phi\phi}$ and they are given by
\begin{align*}
    R_{\theta\theta} &= g^{\phi\phi}R_{\theta\phi\theta\phi}
    = \frac{1}{r^2\sin^2\theta} \cdot r^2\sin^2\theta = 1\\
    R_{\phi\phi} &= g^{\theta\theta}R_{\theta\phi\theta\phi}
    = \frac{1}{r^2} \cdot r^2\sin^2\theta = \sin^2\theta
\end{align*}
Hence the curvature scalar is
\begin{align*}
    R &= g^{\theta\theta}R_{\theta\theta} + g^{\phi\phi}R_{\phi\phi}
    = \frac{1}{r^2}\cdot 1 + \frac{1}{r^2\sin^2\theta}\cdot\sin^2\theta
    = \frac{2}{r^2}
\end{align*}
\end{proof}

\cleardoublepage
\begin{proof}{\textbf{P19.1}}\\
Let us consider the identity (19.4d) but changing names $\alpha\to\beta$ and
$\beta\to\alpha$
\begin{align*}
    R_{\beta\alpha\mu\nu} + R_{\beta\nu\alpha\mu}
    + R_{\beta\mu\nu\alpha} = 0
\end{align*}
Multiplying the whole equation by $g^{\alpha\beta}$ we get that
\begin{align*}
    g^{\alpha\beta}R_{\beta\alpha\mu\nu} + g^{\alpha\beta}R_{\beta\nu\alpha\mu}
    + g^{\alpha\beta}R_{\beta\mu\nu\alpha} &= 0\\
    {R^\alpha}_{\alpha\mu\nu} + {R^{\alpha}}_{\nu\alpha\mu}
    + {R^\alpha}_{\mu\nu\alpha} &= 0\\
    {R^\alpha}_{\alpha\mu\nu} + {R^{\alpha}}_{\mu\alpha\nu}
    - {R^\alpha}_{\mu\alpha\nu} &= 0\\
    {R^\alpha}_{\alpha\mu\nu} &= 0
\end{align*}
Where we used that ${R^{\alpha}}_{\nu\alpha\mu} = +{R^{\alpha}}_{\mu\alpha\nu}$
and that ${R^\alpha}_{\mu\nu\alpha} = -{R^\alpha}_{\mu\alpha\nu}$
\end{proof}

\cleardoublepage
\begin{proof}{\textbf{P19.2}}\\
From the equations $R_{\alpha\beta\mu\nu} = -R_{\alpha\beta\nu\mu}$
and $R_{\alpha\beta\mu\nu} = -R_{\beta\alpha\mu\nu}$ a component of the 
Riemann tensor can be nonzero only if $\alpha \neq \beta$ and $\mu\neq\nu$.
\\
Let's refer to the coordinates as $0, 1$ and $2$ then the only nonzero index
value pairs for $\alpha\beta$ or $\mu\nu$ are $01, 02, 12$.
\\
Let's arrange the nonzero components as in the chart of BOX 19.3
\begin{align*}
    \begin{matrix}
                          &\mu\nu\to & 01       & 02       & 12       \\
    \alpha\beta\downarrow &          &          &          &          \\[7pt]
                       01 &          & R_{0101} & R_{0102} & R_{0112} \\
                       02 &          & R_{0201} & R_{0202} & R_{0212} \\
                       12 &          & R_{1201} & R_{1202} & R_{1212} \\
    \end{matrix}
\end{align*}
But given that we also have the identity
$R_{\alpha\beta\mu\nu} = R_{\mu\nu\alpha\beta}$, the only the elements in the
upper (or down) triangular part of the table are independent.
\\
Therefore in a three dimensional space there are only 6 independent components
of the Riemann tensor.
\end{proof}

\cleardoublepage
\begin{proof}{\textbf{P19.3}}\\
Since must be that $\alpha \neq \beta$ and $\mu \neq \nu$ then,
we want first, to determine the number of possible two-index combinations.
We see that
\begin{align*}
\begin{matrix}
    01 & 02 & 03 & 04 & ... & 0n &   & \text{($n-1$ elements)}\\
       & 12 & 13 & 14 & ... & 1n &   & \text{($n-2$ elements)}\\
       &    & 23 & 24 & ... & 2n &   & \text{($n-3$ elements)}\\
     &  &  & \vdots &  & \vdots &   & \\
       &    &    &    &     & (n-1)n &   & \text{($1$ element)}\\
\end{matrix}
\end{align*}
Then the number of possible two-index combinations is
\begin{align*}
    (n-1) + (n-2) + ... + 1 = \frac{(n - 1)((n-1) + 1)}{2} = \frac{n(n-1)}{2}
\end{align*}
From the hint, we know that the number of off-diagonal elements of an $m \times m$
matrix is $m^2 - m$ so in this case, we are interested in the number of elements
of the upper triangular side of the matrix i.e. we want $(m^2 - m)/2 + m$ where
$m = n(n-1)/2$ hence
\begin{align*}
    \frac{(n(n-1)/2)^2 - n(n-1)/2}{2} + \frac{n(n-1)}{2}
    = \frac{n}{8}(n^3 - 2n^2 + 3n -2) \qquad (*)
\end{align*}
We may write equation 19.4d as
\begin{align*}
    R_{\alpha\beta\mu\nu} = -R_{\alpha\nu\beta\mu} -R_{\alpha\mu\nu\beta}
\end{align*}
So, if the LHS has any repeating indices then the equation becomes the equations
19.4b or 19.4c. For example let us take $R_{0112}$ then we get that
\begin{align*}
    R_{0112} = -R_{0211} -R_{0121} = 0 -R_{0121} = -R_{0121}
\end{align*}
Then, to get a new constraint equation we need to get all four indices
different where order doesn't matter.
\\
The number of elements that satisfy this is $\binom{n}{4}$ i.e.
\begin{align*}
    \binom{n}{4} = \frac{n(n-1)(n-2)(n-3)}{4!}
    = \frac{n(n-1)(n-2)(n-3)}{24}
\end{align*}
So we must subtract all these elements from the ones we counted in $(*)$
wince they are not independent, hence 
\begin{align*}
    \bigg(\frac{n^4}{8} - \frac{n^3}{4} + \frac{3n^2}{8} -\frac{n}{4}\bigg)
    - \bigg(\frac{n^4}{24} - \frac{n^3}{4} + \frac{11n^2}{24} -\frac{n}{4}\bigg)
    = \frac{n^4}{12} - \frac{n^2}{12} = \frac{n^2}{12}(n^2 - 1)
\end{align*}
\end{proof}

\cleardoublepage
\begin{proof}{\textbf{P19.5}}\\
Let us consider the "semilog" coordinates $p,q$ whose metric is
\begin{align*}
    ds^2 = dp^2 + \frac{dq^2}{(bq)^2}
\end{align*}
Let us write the geodesic equation for $\mu = p$ 
\begin{align*}
    \dv{\tau}(g_{p\nu}\dv{x^{\nu}}{\tau})
    - \frac{1}{2}\partial_{p}g_{\alpha\beta}\dv{x^\alpha}{\tau}\dv{x^\beta}{\tau}
    &= 0\\
    \dv[2]{p}{\tau} - 0 &= 0
\end{align*}
So all Christoffel symbols with $p$ as a superscript are 0.
\\
For $\mu = q$ we have that
\begin{align*}
    \dv{\tau}(g_{q\nu}\dv{x^{\nu}}{\tau})
    - \frac{1}{2}\partial_{q}g_{\alpha\beta}\dv{x^\alpha}{\tau}\dv{x^\beta}{\tau}
    &= 0\\
    \dv{\tau}(\frac{1}{b^2q^2}\dv{q}{\tau})
    + \frac{1}{2}\frac{2}{b^2q^3}\bigg(\dv{q}{\tau}\bigg)^2
    &= 0\\
    -\frac{2}{b^2q^3}\bigg(\dv{q}{\tau}\bigg)^2
    + \frac{1}{b^2q^2}\dv[2]{q}{\tau}
    + \frac{1}{b^2q^3}\bigg(\dv{q}{\tau}\bigg)^2
    &= 0\\
    \dv[2]{q}{\tau}
    - \frac{1}{q}\bigg(\dv{q}{\tau}\bigg)^2
    &= 0
\end{align*}
Then the only non-zero Christoffel symbols with $q$ as a superscript is
\begin{align*}
    \Gamma^q_{qq} = -\frac{1}{q}
\end{align*}
Since we are in 2 dimensions, then by BOX 19.3 we know there is only 1
independent component of the Riemann tensor that satisfies
\begin{align*}
    R_{pqpq} = -R_{pqqp} = R_{qpqp} = -R_{qppq}
\end{align*}
So we compute only $R_{pqpq}$ as follows
\begin{align*}
    R_{pqpq} &= g_{p\alpha}{R^\alpha}_{qpq}\\
    &= g_{pp}{R^p}_{qpq} + 0\\
    &= \partial_p\Gamma^p_{qq} - \partial_q\Gamma^p_{qp}
    + \Gamma^p_{p\gamma} \Gamma^\gamma_{qq}
    - \Gamma^p_{q\sigma} \Gamma^\sigma_{qp}\\
    &= 0 - 0 + 0 - 0 = 0
\end{align*}
So all Riemann components are 0.\\
Given that all Riemann tensor components are 0 then, all the components of
the Ricci tensor are 0
\begin{align*}
    R_{\beta\nu} = {R^\alpha}_{\beta\alpha\nu} = 0
\end{align*}
In the same way, the curvature scalar is also 0
\begin{align*}
    R = g^{\beta\nu}R_{\beta\nu} = 0
\end{align*}
Therefore this metric describes a flat two-dimensional space.
\end{proof}

\cleardoublepage
\begin{proof}{\textbf{P19.6}}
Let us consider the following metric
$$ds^2 = dp^2 + e^{2p/p_0}dq^2$$
\begin{itemize}
\item [\textbf{a.}] Let us write the geodesic equation for $\mu = p$
\begin{align*}
    \dv{\tau}(g_{p\nu}\dv{x^{\nu}}{\tau})
    - \frac{1}{2}\partial_{p}g_{\alpha\beta}\dv{x^\alpha}{\tau}\dv{x^\beta}{\tau}
    &= 0\\
    \dv[2]{p}{\tau}
    - \frac{1}{2}\frac{2}{p_0}e^{2p/p_0}\bigg(\dv{q}{\tau}\bigg)^2
    &= 0\\
    \dv[2]{p}{\tau}
    - \frac{e^{2p/p_0}}{p_0}\bigg(\dv{q}{\tau}\bigg)^2
    &= 0
\end{align*}
Then the only non-zero Christoffel symbol with $p$ as a superscript is
\begin{align*}
    \Gamma^p_{qq} = - \frac{e^{2p/p_0}}{p_0}
\end{align*}
For $\mu = q$ we get that
\begin{align*}
    \dv{\tau}(g_{q\nu}\dv{x^{\nu}}{\tau})
    - \frac{1}{2}\partial_{q}g_{\alpha\beta}\dv{x^\alpha}{\tau}\dv{x^\beta}{\tau}
    &= 0\\
    \dv{\tau}(e^{2p/p_0}\dv{q}{\tau}) - 0
    &= 0\\
    \frac{2e^{2p/p_0}}{p_0}\dv{p}{\tau}\dv{q}{\tau} + e^{2p/p_0}\dv[2]{q}{\tau}
    &= 0\\
    \dv[2]{q}{\tau} + \frac{2}{p_0}\dv{p}{\tau}\dv{q}{\tau}
    &= 0
\end{align*}
Then the only non-zero Christoffel symbols with $q$ as a superscript are
\begin{align*}
    \Gamma^q_{pq} = \Gamma^q_{qp} = \frac{1}{p_0}
\end{align*}

\item [\textbf{b.}] Let us compute the only independent Riemann tensor component
as follows
\begin{align*}
    R_{pqpq} &= g_{p\alpha}{R^\alpha}_{qpq}\\
    &= g_{pp}{R^p}_{qpq} + 0\\
    &= \partial_p\Gamma^p_{qq} - \partial_q\Gamma^p_{qp}
    + \Gamma^p_{p\gamma} \Gamma^\gamma_{qq}
    - \Gamma^p_{q\sigma} \Gamma^\sigma_{qp}\\
    &= - \frac{2e^{2p/p_0}}{p_0^2} - 0  + 0
    - \bigg(- \frac{e^{2p/p_0}}{p_0}\bigg)\bigg(\frac{1}{p_0}\bigg)\\
    &= - \frac{2e^{2p/p_0}}{p_0^2} + \frac{e^{2p/p_0}}{p_0^2}\\
    &= -\frac{e^{2p/p_0}}{p_0^2}
\end{align*}
Since there is at least one non-zero component of the Riemann tensor,
we can say that the 2D space is curved.
\end{itemize}
\end{proof}

\cleardoublepage
\begin{proof}{\textbf{P19.7}}
Let us consider the metric
\begin{align*}
    ds^2 = dr^2 + r^2d\theta^2 + r^2\sin^2\theta d\phi^2
\end{align*}
\begin{itemize}
\item [\textbf{a.}] Let us consider the geodesic equation when $\mu = r$ then
\begin{align*}
    \dv{\tau}(g_{r\nu}\dv{x^{\nu}}{\tau})
    - \frac{1}{2}\partial_{r}g_{\alpha\beta}\dv{x^\alpha}{\tau}\dv{x^\beta}{\tau}
    &= 0\\
    \dv[2]{r}{\tau}
    - \frac{1}{2}\bigg[
        2r\bigg(\dv{\theta}{\tau}\bigg)^2
        + 2r\sin^2\theta\bigg(\dv{\phi}{\tau}\bigg)^2
    \bigg]
    &= 0\\
    \dv[2]{r}{\tau}
    - r\bigg(\dv{\theta}{\tau}\bigg)^2
    - r\sin^2\theta\bigg(\dv{\phi}{\tau}\bigg)^2
    &= 0
\end{align*}
Then the only non-zero Christoffel symbols with $r$ as a superscript are
\begin{align*}
    \Gamma^r_{\theta\theta} = -r
    \qquad
    \Gamma^r_{\phi\phi} = -r\sin^2\theta
\end{align*}
If we let $\mu = \theta$ we get that
\begin{align*}
    \dv{\tau}(g_{\theta\nu}\dv{x^{\nu}}{\tau})
    - \frac{1}{2}\partial_{\theta}
    g_{\alpha\beta}\dv{x^\alpha}{\tau}\dv{x^\beta}{\tau}
    &= 0\\
    \dv{\tau}(r^2\dv{\theta}{\tau})
    - r^2\sin\theta\cos\theta\bigg(\dv{\phi}{\tau}\bigg)^2
    &= 0\\
    2r\dv{r}{\tau}\dv{\theta}{\tau} + r^2\dv[2]{\theta}{\tau} 
    - r^2\sin\theta\cos\theta\bigg(\dv{\phi}{\tau}\bigg)^2
    &= 0\\
    \dv[2]{\theta}{\tau} + \frac{2}{r}\dv{r}{\tau}\dv{\theta}{\tau}
    - \sin\theta\cos\theta\bigg(\dv{\phi}{\tau}\bigg)^2
    &= 0
\end{align*}
Then the only non-zero Christoffel symbols with $\theta$ as a superscript are
\begin{align*}
    \Gamma^\theta_{r\theta} = \Gamma^\theta_{\theta r} = \frac{1}{r}
    \qquad
    \Gamma^\theta_{\phi\phi} = - \sin\theta\cos\theta
\end{align*}
We already knew the last one from the BOX 19.6.
\\
Finally, if we let $\mu = \phi$ we get that
\begin{align*}
    \dv{\tau}(g_{\phi\nu}\dv{x^{\nu}}{\tau})
    - \frac{1}{2}\partial_{\phi}
    g_{\alpha\beta}\dv{x^\alpha}{\tau}\dv{x^\beta}{\tau}
    &= 0\\
    \dv{\tau}(r^2\sin^2\theta\dv{\phi}{\tau})
    - 0 &= 0\\
    2r\sin^2\theta\dv{r}{\tau}\dv{\phi}{\tau}
    + r^2\bigg(2\sin\theta\cos\theta\dv{\theta}{\tau}\dv{\phi}{\tau}
    + \sin^2\theta\dv[2]{\phi}{\tau}\bigg)
    &= 0\\
    \dv[2]{\phi}{\tau}
    + \frac{2}{r}\dv{r}{\tau}\dv{\phi}{\tau}
    + 2\frac{\cos\theta}{\sin\theta}\dv{\theta}{\tau}\dv{\phi}{\tau}
    &= 0
\end{align*}
So the only non-zero Christoffel symbols with $\phi$ as a superscript are
\begin{align*}
    \Gamma^\phi_{r\phi} = \Gamma^\phi_{\phi r} = \frac{1}{r}
    \qquad
    \Gamma^\phi_{\theta\phi} = \Gamma^\phi_{\phi\theta} = \cot\theta
\end{align*}
Again, we already knew the last one from the BOX 19.6.

\item[\textbf{b.}]
We know that in this case, the Riemann tensor component
$R_{\theta\phi\theta\phi}$ is given by
\begin{align*}
    R_{\theta\phi\theta\phi} =
    r^2(\partial_\theta\Gamma^\theta_{\phi\phi}
    + \Gamma^\theta_{\theta\gamma}\Gamma^\gamma_{\phi\phi}
    - \Gamma^\theta_{\phi\sigma}\Gamma^\sigma_{\phi\theta})
\end{align*}
We cannot remove the middle term since
$\Gamma^\theta_{\theta r}$ and $\Gamma^r_{\phi\phi}$ are non-zero, then we get
that
\begin{align*}
    R_{\theta\phi\theta\phi} &=
    r^2(\partial_\theta\Gamma^\theta_{\phi\phi}
    - \Gamma^\theta_{\phi\phi}\Gamma^\phi_{\phi\theta})
    + r^2\Gamma^\theta_{\theta r}\Gamma^r_{\phi\phi}\\
    &= r^2((\sin^2\theta - \cos^2\theta)
    + (\sin\theta\cos\theta) \cdot \cot\theta)
    - r^2\bigg(\frac{1}{r}\cdot r\sin^2\theta\bigg)\\
    &= r^2(\sin^2\theta - \cos^2\theta + \cos^2\theta)
    - r^2\sin^2\theta\\
    &= r^2\sin^2\theta - r^2\sin^2\theta\\
    &= 0
\end{align*}
Therefore we get an additional term $r^2\sin^2\theta$ which cancels the 
$r^2\sin^2\theta$ we had when we considered $r$ as a constant.
\end{itemize}
\end{proof}

\cleardoublepage
\begin{proof}{\textbf{P19.8}}\\
\begin{itemize}
\item [\textbf{a.}]
In any coordinate basis we note that
\begin{align*}
    \bm{R}
    &= R_{\zeta\gamma\delta\sigma}\bm{e}^\zeta\bm{e}^\gamma\bm{e}^\delta\bm{e}^\sigma \\
    &= g_{\zeta\lambda}{R^\lambda}_{\gamma\delta\sigma}\bm{e}^\zeta\bm{e}^\gamma\bm{e}^\delta\bm{e}^\sigma 
\end{align*}
So
\begin{align*}
    (R_{\rho\beta\mu\nu})_{obs}
    &= (\bm{o}_\rho)^\zeta (\bm{o}_\beta)^\gamma (\bm{o}_\mu)^\delta (\bm{o}_\nu)^\sigma
    g_{\zeta\lambda}{R^\lambda}_{\gamma\delta\sigma} 
\end{align*}
Now, raising the index $\alpha$ we get that
\begin{align*}
    ({R^\alpha}_{\beta\mu\nu})_{obs}
    &= \eta^{\alpha\rho}(R_{\rho\beta\mu\nu})_{obs}\\
    &= \eta^{\alpha\rho}
    (\bm{o}_\rho)^\zeta (\bm{o}_\beta)^\gamma (\bm{o}_\mu)^\delta (\bm{o}_\nu)^\sigma
    g_{\zeta\lambda}{R^\lambda}_{\gamma\delta\sigma} 
\end{align*}

\item [\textbf{b.}] Let us consider the Riemann tensor component
$({R^t}_{ztz})_{obs}$ then, we see that
\begin{align*}
    ({R^t}_{ztz})_{obs}
    &= \eta^{t\rho}
    (\bm{o}_\rho)^\zeta (\bm{o}_z)^\gamma (\bm{o}_t)^\delta (\bm{o}_z)^\sigma
    g_{\zeta\lambda}{R^\lambda}_{\gamma\delta\sigma} \\
    &= -(\bm{o}_t)^\zeta (\bm{o}_z)^\gamma (\bm{o}_t)^\delta (\bm{o}_z)^\sigma
    g_{\zeta\lambda}{R^\lambda}_{\gamma\delta\sigma} 
\end{align*}
Also, we know that for an observer falling radially we have that
\begin{align*}
    (\bm{o}_t)^{\mu} = \begin{bmatrix}
        \left(1 - \frac{2GM}{r}\right)^{-1}\\[7pt]
        -\sqrt{\frac{2GM}{r}}\\[7pt]
        0\\
        0
    \end{bmatrix}
    \qquad
    (\bm{o}_z)^{\mu} = \begin{bmatrix}
        \frac{-\sqrt{2GM/r}}{1 - 2GM/r}\\[7pt]
        1\\
        0\\
        0
    \end{bmatrix}
\end{align*}
Hence
% Because of the symmetries of the Riemann tensor we only need to evaluate the
% component ${R^t}_{rtr}$, hence
\begin{align*}
    ({R^t}_{ztz})_{obs}
    &= -(
    (\bm{o}_t)^t (\bm{o}_z)^r (\bm{o}_t)^t (\bm{o}_z)^r g_{tt}{R^t}_{rtr} 
    + (\bm{o}_t)^t (\bm{o}_z)^r (\bm{o}_t)^r (\bm{o}_z)^t g_{tt}{R^t}_{rrt}\\
    &\quad
    + (\bm{o}_t)^r (\bm{o}_z)^t (\bm{o}_t)^r (\bm{o}_z)^t g_{rr}{R^r}_{trt}
    + (\bm{o}_t)^r (\bm{o}_z)^t (\bm{o}_t)^t (\bm{o}_z)^r g_{rr}{R^r}_{ttr}
    )\\[7pt]
    &= -(
    (\bm{o}_t)^t (\bm{o}_z)^r (\bm{o}_t)^t (\bm{o}_z)^r R_{trtr} 
    + (\bm{o}_t)^t (\bm{o}_z)^r (\bm{o}_t)^r (\bm{o}_z)^t R_{trrt}\\
    &\quad
    + (\bm{o}_t)^r (\bm{o}_z)^t (\bm{o}_t)^r (\bm{o}_z)^t R_{rtrt}
    + (\bm{o}_t)^r (\bm{o}_z)^t (\bm{o}_t)^t (\bm{o}_z)^r R_{rttr}
    )\\[7pt]
    &= -(
    (\bm{o}_t)^t (\bm{o}_z)^r (\bm{o}_t)^t (\bm{o}_z)^r R_{trtr} 
    - (\bm{o}_t)^t (\bm{o}_z)^r (\bm{o}_t)^r (\bm{o}_z)^t R_{trtr}\\
    &\quad
    + (\bm{o}_t)^r (\bm{o}_z)^t (\bm{o}_t)^r (\bm{o}_z)^t R_{trtr}
    - (\bm{o}_t)^r (\bm{o}_z)^t (\bm{o}_t)^t (\bm{o}_z)^r R_{trtr}
    )\\[7pt]
    &= -(
    (\bm{o}_t)^t (\bm{o}_z)^r (\bm{o}_t)^t (\bm{o}_z)^r
    + (\bm{o}_t)^r (\bm{o}_z)^t (\bm{o}_t)^r (\bm{o}_z)^t\\
    &\quad- 2(\bm{o}_t)^t (\bm{o}_z)^r (\bm{o}_t)^r (\bm{o}_z)^t
    ) R_{trtr}\\[7pt]
    &= -(
    (\bm{o}_t)^t (\bm{o}_t)^t
    + (\bm{o}_t)^r (\bm{o}_z)^t (\bm{o}_t)^r (\bm{o}_z)^t
    - 2(\bm{o}_t)^t (\bm{o}_t)^r (\bm{o}_z)^t
    ) R_{trtr}
\end{align*}
Where we used that $(\bm{o}_z)^r = 1$.\\
Now, we replace all the other values as follows
\begin{align*}
    ({R^t}_{ztz})_{obs}
    &= -\bigg(\bigg(1 -\frac{2GM}{r}\bigg)^{-2} 
    + \frac{2GM}{r}\frac{2GM}{r}\bigg(1 -\frac{2GM}{r}\bigg)^{-2}\\
    &\quad- 2\bigg(1 -\frac{2GM}{r}\bigg)^{-1}\frac{2GM}{r}
    \bigg(1 -\frac{2GM}{r}\bigg)^{-1}\bigg) R_{trtr}\\[7pt]
    &= -\bigg(
    1 - 2\frac{2GM}{r} + \frac{(2GM)^2}{r^2}
    \bigg)\bigg(1 -\frac{2GM}{r}\bigg)^{-2} R_{trtr}\\[7pt]
    &= -\bigg(1 - \frac{2GM}{r}\bigg)^2
    \bigg(1 -\frac{2GM}{r}\bigg)^{-2} R_{trtr}\\[7pt]
    &= - g_{tt} {R^t}_{rtr}\\[7pt]
    &= \bigg(1 - \frac{2GM}{r}\bigg)
    \frac{2GM}{r^3}\bigg(1 - \frac{2GM}{r}\bigg)^{-1}\\
    &= \frac{2GM}{r^3}
\end{align*}
\end{itemize}
\end{proof}

\cleardoublepage
\begin{proof}{\textbf{P19.9}}\\
From problem 18.7 we know that the extremities will experience a force given by
\begin{align*}
    m\dv[2]{n^z}{t} = -m{R^z}_{tzt} n^z
\end{align*}
Because they are in the same locally inertial frame as the center of mass.
\\
Also, we note that
\begin{align*}
    {R^z}_{tzt} = g^{zz}R_{ztzt} = -g^{zz}R_{tzzt} =
    g^{zz}R_{tztz} = g^{zz}g^{tt}{R^t}_{ztz} = -{R^t}_{ztz}
\end{align*}
Where we used that $g^{zz} = g^{rr}$ in the Schwarzschild metric. So
\begin{align*}
    m\dv[2]{n^z}{t} &= m{R^t}_{ztz} n^z
    = m\frac{2GM}{r^3} n^z
\end{align*}
If we take the force to be $F_{max}$ and we solve for $r = R$ we get that 
\begin{align*}
    R = \bigg(\frac{2GMm}{F_{max}}n^z\bigg)^{1/3}
\end{align*}
So, at $r = R$ the force gets to $F_{max}$ where we start feeling serious pain.
\\
Now, from equation (14.4) we see that the proper time measured by a particle 
falling from rest at $r = R$ to $r = 0$ is
\begin{align*}
    \Delta \tau &= \frac{\pi}{2}\sqrt{\frac{R^3}{2GM}}
    = \frac{\pi}{2}\sqrt{\frac{mn^z}{F_{max}}}
\end{align*}
Assuming that $F_{max} = 200~N$, $n^z = 1~m$ and $m = 30~kg$ then this time is
\begin{align*}
    \Delta \tau 
    &= \frac{\pi}{2}\sqrt{\frac{30}{200}} = 0.60~s
\end{align*}
Since the pain impulses travel at about $1~m/s$ along the nerves, then the
person would not feel any pain since it would take 1s to arrive the pain signal
to the brain from the extremities, but we fall faster than that into $r = 0$. 
\end{proof}
\end{document}